%%% FIELD proof-local-minimax-attainment
Let
\[
L_{b,R}:=T_b(s_0)\cap\{h\in\mathbb R^d:\|h\|_2\le R\},
\qquad b\in\mathbb B(s_0).
\]
The hypothesis \(\Gamma(s_0)=\{\tau_0\}\), together with the definition of \(\Gamma(s_0)\), implies that \(\mathbb B(s_0)\) is nonempty.  It is finite because it is a subset of the finite set \(\mathbb B_p\).  Each \(T_b(s_0)\) is closed by the declared type of the tangent cone, so \(L_{b,R}\) is a closed subset of the compact Euclidean ball of radius \(R\), and hence is compact.  By the definition of the branch-local target, \(\eta_b\) is linear and therefore continuous.  Consequently
\[
K_R
=\bigcup_{b\in\mathbb B(s_0)}\eta_b(L_{b,R})
\]
is a nonempty finite union of compact sets and is compact.  In particular,
\[
D_R:=\operatorname{diam}(K_R)<\infty.
\]

We first justify the restriction to actions in \(\mathcal H_R\).  For a closed set \(C\subseteq\mathbb R^3\), define
\[
T(C):=C\cap K_R.
\]
This map is Effros measurable.  Indeed, let \(O=K_R\cap G\) be relatively open in \(K_R\), where \(G\subseteq\mathbb R^3\) is open, and define
\[
A_m:=\{x\in K_R:\operatorname{dist}(x,G^c)\ge m^{-1}\},
\qquad m\ge1,
\]
with \(A_m=K_R\) if \(G^c=\varnothing\).  Each \(A_m\) is compact and \(O=\bigcup_{m\ge1}A_m\).  If \(A\) is compact, closedness of \(C\) and compactness of \(A\) give
\[
\{C:C\cap A\ne\varnothing\}
=\bigcap_{\ell\ge1}\{C:C\cap A^{1/\ell}\ne\varnothing\},
\]
where \(A^{1/\ell}\) is the open \(1/\ell\)-neighborhood of \(A\).  The events on the right are Effros hit events.  Thus compact-hit events are measurable, and
\[
\{C:T(C)\cap O\ne\varnothing\}
=\bigcup_{m\ge1}\{C:C\cap A_m\ne\varnothing\}
\]
is measurable.  The inverse image of the isolated empty element is
\[
\{C:T(C)=\varnothing\}=\{C:C\cap K_R=\varnothing\},
\]
which is also measurable by the compact-hit identity.  These events generate the Effros sigma-field of \(\mathcal H_R\), proving the assertion.

For every \((b,h)\in\Lambda_R\), the definition of \(K_R\) gives \(\eta_b(h)\in K_R\).  Hence, pointwise in every closed \(C\),
\[
{\bf1}\{\eta_b(h)\in T(C)\}
={\bf1}\{\eta_b(h)\in C\},
\qquad
\operatorname{diam}(T(C))\le\operatorname{diam}(C).
\]
Thus truncation preserves all coverage probabilities and cannot increase any expected-diameter risk.  By \(\text{lem:hyperspace-compactness}\), \(\mathcal H_R\) is compact and metrizable, hence standard Borel, and \(C\mapsto\operatorname{diam}(C)\) is continuous on \(\mathcal H_R\).  It is bounded there by \(D_R\).

We next verify that all Gaussian shifts admit the common domination used below.  Write \(\Sigma_{0,e}\) for the positive-definite matrix in environment \(e\) represented by \(s_0\).  For a nonzero half-vector coordinate \(c\), let \(A_c\ne0\) be the unique symmetric matrix satisfying
\[
c^{\mathsf T}\operatorname{vech}(zz^{\mathsf T})=z^{\mathsf T}A_cz
\quad\text{for every }z\in\mathbb R^p.
\]
A direct expansion using the displayed coordinates of \(V(\Sigma)\) yields
\[
c^{\mathsf T}V(\Sigma)c
=2\operatorname{tr}\!\left[
  \{\Sigma^{1/2}A_c\Sigma^{1/2}\}^2
 \right]>0
\]
whenever \(\Sigma\) is positive definite: the symmetric matrix \(\Sigma^{1/2}A_c\Sigma^{1/2}\) is nonzero, and the trace of its square is the sum of the squares of its eigenvalues.  Since \(\lambda\in(0,1)\), the block-diagonal definition of \(V_0\) therefore implies \(V_0\succ0\).

Let \(\mu_0=\mathcal N(0,V_0)\).  For each \(h\in\mathbb R^d\), \(\mathcal N(h,V_0)\) is equivalent to \(\mu_0\), with continuous strictly positive likelihood ratio
\[
w_h(y):=\exp\!\left{
h^{\mathsf T}V_0^{-1}y-	frac12h^{\mathsf T}V_0^{-1}h
\right},
\qquad
\varphi_h(y)\,dy=w_h(y)\,\mu_0(dy),
\qquad
\int w_h\,d\mu_0=1.
\]

Let \(\delta\) be an \(\mathcal H_R\)-valued randomized confidence rule.  Measurability of \(\delta\) and independence and uniformity of \(U\) define the probability kernel
\[
Q(A\mid y):=\int_0^1{\bf1}\{\delta(y,u)\in A\}\,du,
\qquad A\in\mathcal B(\mathcal H_R).
\]
Conversely, because \(\mathcal H_R\) is standard Borel, \(\text{lem:kernel-uniform-randomization}\) realizes every probability kernel \(Q\) from \(\mathbb R^d\) to \(\mathcal H_R\) as the conditional law of a jointly measurable function of \((y,U)\).  For \(k=(b,h)\in\Lambda_R\), integrating the conditional law of the randomized action and then applying the preceding change of measure shows that the rule's coverage probability is
\[
\int_{\mathbb R^d}\!\int_{\mathcal H_R}
 {\bf1}\{\eta_b(h)\in C\}\,Q(dC\mid y)\varphi_h(y)\,dy,
\]
and its expected diameter is
\[
\int_{\mathbb R^d}\!\int_{\mathcal H_R}
 \operatorname{diam}(C)\,Q(dC\mid y)\varphi_h(y)\,dy.
\]
Therefore the feasible kernels, together with an epigraph coordinate \(t\), are exactly the variables and constraints in the displayed semi-infinite linear program, and the infimum of that program equals \(v_R\).  The deterministic rule \(C=K_R\) has coverage one and risk \(D_R\).  Hence
\[
0\le v_R\le D_R<\infty.
\]

It remains to prove attainment.  For every integer \(j\ge1\), choose a feasible kernel \(Q_j\) whose worst-case expected diameter is at most \(v_R+j^{-1}\), and put
\[
\zeta_j(dy,dC):=\mu_0(dy)Q_j(dC\mid y)
\]
on \(\mathbb R^d\times\mathcal H_R\).  The first marginal of every \(\zeta_j\) is \(\mu_0\).  Given \(\epsilon>0\), choose a compact Euclidean ball \(L_\epsilon\) such that \(\mu_0(L_\epsilon)>1-\epsilon\).  Then the compact set \(L_\epsilon\times\mathcal H_R\) has \(\zeta_j\)-probability greater than \(1-\epsilon\) for every \(j\), so \((\zeta_j)_{j\ge1}\) is tight.  By \(\text{lem:prokhorov-tightness}\), a subsequence, not relabeled, converges weakly to a probability measure \(\zeta\).

For every bounded continuous \(g:\mathbb R^d\to\mathbb R\), weak convergence gives
\[
\int g(y)\,\zeta(dy,dC)
=\lim_{j\to\infty}\int g(y)\,\zeta_j(dy,dC)
=\int g(y)\,\mu_0(dy).
\]
Thus \(\zeta\) has first marginal \(\mu_0\).  Since both factors are standard Borel spaces, \(\text{lem:standard-borel-disintegration}\) supplies a probability kernel \(Q\) such that
\[
\zeta(dy,dC)=\mu_0(dy)Q(dC\mid y).
\]

Fix \(k=(b,h)\in\Lambda_R\), and define
\[
F_k:=\{(y,C):\eta_b(h)\in C\}.
\]
Membership of a fixed point is closed under Hausdorff convergence of nonempty compact sets, and the empty set is isolated in \(\mathcal H_R\); hence \(F_k\) is closed.  For \(L>0\), let \(w_h^L:=w_h\wedge L\).  The function \(w_h^L{\bf1}_{F_k}\) is bounded and upper semicontinuous.  By \(\text{lem:portmanteau-usc}\), feasibility of \(Q_j\), and the fact that every \(\zeta_j\) has first marginal \(\mu_0\),
\[
\begin{aligned}
\int w_h^L{\bf1}_{F_k}\,d\zeta
&\ge\limsup_{j\to\infty}\int w_h^L{\bf1}_{F_k}\,d\zeta_j\\
&\ge 1-\gamma-\int (w_h-w_h^L)\,d\mu_0.
\end{aligned}
\]
Because \(w_h\) is nonnegative and \(\mu_0\)-integrable,
\[
\int (w_h-w_h^L)\,d\mu_0\longrightarrow0
\quad\text{as }L\to\infty.
\]
Monotone convergence on the left therefore gives
\[
\int w_h{\bf1}_{F_k}\,d\zeta\ge1-\gamma.
\]
This is the coverage constraint at \(k\); since \(k\) was arbitrary, \(Q\) is uniformly coverage feasible.

For the risk constraint, the function
\[
(y,C)\longmapsto w_h^L(y)\operatorname{diam}(C)
\]
is bounded and continuous.  Weak convergence yields
\[
\int w_h^L(y)\operatorname{diam}(C)\,d\zeta
=\lim_{j\to\infty}
 \int w_h^L(y)\operatorname{diam}(C)\,d\zeta_j.
\]
For both \(\zeta_j\) and \(\zeta\), whose first marginal is \(\mu_0\), the error from replacing \(w_h^L\) by \(w_h\) is bounded by
\[
D_R\int(w_h-w_h^L)\,d\mu_0\longrightarrow0.
\]
It follows that, for this fixed \(k\),
\[
\begin{aligned}
\int w_h(y)\operatorname{diam}(C)\,d\zeta
&=\lim_{j\to\infty}
 \int w_h(y)\operatorname{diam}(C)\,d\zeta_j\\
&\le\limsup_{j\to\infty}(v_R+j^{-1})=v_R.
\end{aligned}
\]
Taking the supremum over \(k\in\Lambda_R\) proves that the feasible kernel \(Q\) has worst-case risk at most \(v_R\).  The definition of \(v_R\) gives the reverse inequality, so this risk equals \(v_R\).  Thus the infimum is attained and the displayed semi-infinite linear program has value \(v_R\).

Every measurable uniform-randomizer realization of a minimizing kernel has the just-computed coverage and risk and is therefore attaining.  Conversely, the kernel induced by an attaining \(\mathcal H_R\)-valued rule is feasible and has objective \(v_R\), so it is minimizing.  Since every \(\mathcal N(h,V_0)\) is equivalent to \(\mu_0\), modifications on \(\mu_0\)-null sets are null under every experiment law, which gives the asserted characterization up to null sets.  If an unrestricted attaining rule is given, the measurable truncation \(T\circ\delta\) remains feasible and has risk at most \(v_R\).  Minimality forces its risk to equal \(v_R\), so its induced kernel is minimizing and has a uniform-randomizer representation.  This does not imply that the original, untruncated rule takes values in \(\mathcal H_R\).

Finally, consider any proposed sequence \((\Delta_N)_{N\ge1}\) of finite-sample pullbacks and fix one index \(N_0\).  Altering only \(\Delta_{N_0}\) leaves the sequence unchanged for every \(N>N_0\), so it leaves every local asymptotic limit as \(N\to\infty\) unchanged.  Replacing \(\Delta_{N_0}\) by the constantly empty set makes its coverage at \(N_0\) equal to zero.  Replacing it instead by the constant set \(\mathbb R^3\) makes its diameter at \(N_0\) equal to \(+\infty\).  Hence the local attainment characterization alone entails neither finite-sample coverage nor finite-sample diameter dominance.
